\labsection{8}{Testing whether a cryptocurrency is money}{sec:lab08-crypto}

\begin{labproject}{Apply the three-function test with evidence rather than opinion}
\textbf{Coverage.} Chapter~\ref{ch:cryptocurrencies}.\\
\textbf{Goal.} Measure the properties Chapter~\ref{ch:cryptocurrencies} discusses qualitatively, and reach a defensible verdict on each of money's three functions.

\begin{labmode}
The price series is synthetic, generated from a seeded random walk calibrated to realistic Bitcoin and fiat volatility. Results are reproducible and the magnitudes are representative, but no figure here is a real historical price --- do not quote one as such.

If you can obtain real daily closes from a public source, rerun the identical script against them. The column layout matches.
\end{labmode}

\medskip\noindent\textbf{Part A --- Measure volatility properly.}\par

\begin{lstlisting}[style=mlpython]
import numpy as np
import pandas as pd

px = pd.read_csv("price_series.csv", parse_dates=["date"]).set_index("date")

# Log returns, because they are additive over time and symmetric in
# direction: a +50% move followed by -50% is not a round trip in simple
# returns, but log returns handle it correctly.
returns = np.log(px / px.shift(1)).dropna()

annualised = returns.std() * np.sqrt(365)   # crypto trades every day
print((annualised * 100).round(1))
\end{lstlisting}

On the supplied series:

\begin{lstlisting}[style=mlterminal]
crypto_usd      69.8%
equity_index    18.6%
fiat_eurusd      8.9%
\end{lstlisting}

The crypto series is roughly eight times as volatile as the currency pair, and nearly four times a broad equity index. That gap is the whole of Part B.

\medskip\noindent\textbf{Part B --- Maximum drawdown.}\par
Volatility treats gains and losses alike. A holder does not. Drawdown measures the worst peak-to-trough fall --- the loss someone who bought at the wrong moment actually experienced:

\begin{lstlisting}[style=mlpython]
def max_drawdown(series: pd.Series) -> float:
    running_peak = series.cummax()
    drawdown = series / running_peak - 1
    return drawdown.min()

for col in px.columns:
    print(f"{col:10s} max drawdown {max_drawdown(px[col]):.1%}")
\end{lstlisting}

\begin{lstlisting}[style=mlterminal]
asset             max drawdown   gain to recover
crypto_usd              -65.8%            192.0%
fiat_eurusd             -21.1%             26.8%
equity_index            -21.1%             26.7%
\end{lstlisting}

\begin{workedexample}
A 65.8\% drawdown means an asset bought at the peak fell to 34.2\% of its value.

Recovering requires a gain of $1/0.342 - 1 = 192\%$ --- the asset must nearly triple simply to return to where the holder started.

Compare the fiat row: a 21.1\% fall needs only 26.8\% to recover. The relationship is not linear. Halving the drawdown does far more than halve the recovery burden, which is why drawdown matters more than volatility for a store of value.
\end{workedexample}

\medskip\noindent\textbf{Part C --- The three-function test.}\par
Chapter~\ref{ch:cryptocurrencies} sets out the standard economic test. Judge each function against a measurement, not an intuition:

\begin{mltable}
\begin{tabularx}{\linewidth}{@{}>{\bfseries\leavevmode\color{MLNavy}}p{0.22\linewidth}X X@{}}
\toprule
\rowcolor{MLPaleBlue!92}
Function & Evidence that satisfies it & Measure it with \\
\midrule
Medium of exchange & Accepted widely; settles quickly; fees small relative to typical purchase & Confirmation time; fee as \% of a \$20 purchase \\
Store of value & Purchasing power roughly preserved over the holding period & Annualised volatility; max drawdown \\
Unit of account & Prices are quoted in it natively & Count real goods priced in the currency rather than pegged to fiat \\
\bottomrule
\end{tabularx}
\end{mltable}

\begin{lstlisting}[style=mlpython]
# Medium of exchange: is the fee proportionate to everyday spending?
purchase = 20.00
for fee in [0.15, 1.20, 8.50]:      # low, typical, congested network
    print(f"fee ${fee:5.2f} = {fee / purchase:6.1%} of a $20 purchase")

# Settlement: six confirmations at ~10 minutes per block.
print(f"\ntypical settlement: {6 * 10} minutes")
\end{lstlisting}

\begin{keyidea}
The three functions are not independent, and the dependency runs in one direction.

Volatility undermines store of value directly. But it also undermines unit of account, because nobody prices a menu in a unit that moves 5\% overnight. And it undermines medium of exchange, because a merchant accepting payment converts to fiat immediately to avoid the exposure --- so the currency is a transfer rail, not money.

Chapter~\ref{ch:cryptocurrencies} presents volatility as one risk among several. Your measurements should show it is the constraint the other two functions hang from.
\end{keyidea}

\begin{pitfall}
``Bitcoin has no intrinsic value'' is a weak argument, and using it will cost you marks.

Fiat currency has no intrinsic value either. It is accepted because a state demands taxes in it and enforces it as legal tender. The interesting question is not intrinsic value but what sustains acceptance --- and the honest comparison is between different \emph{sources} of confidence, not between one asset that has backing and one that does not.
\end{pitfall}

\medskip\noindent\textbf{Part D --- The trade-off nobody escapes.}\par
Chapter~\ref{ch:cryptocurrencies} raises user control as a benefit. Price it honestly:

\begin{itemize}
  \item Self-custody means no counterparty can freeze or seize your funds.
  \item It also means a lost private key is a permanent, unrecoverable loss --- with no fraud department and no chargeback.
  \item Custodial exchanges restore recovery and support, and reintroduce exactly the trusted intermediary the design set out to remove.
\end{itemize}

\begin{tryit}
Estimate how much Bitcoin is permanently lost to forgotten keys. Published estimates cluster around 3--4 million coins; find one and cite it.

Then take a position on the design question that follows: is irreversibility a feature or a defect? Argue it separately for a large cross-border settlement and for a consumer buying coffee. A good answer reaches different conclusions for the two cases and explains why.
\end{tryit}
\end{labproject}

\begin{verification}
\begin{itemize}
  \item Returns are computed as log returns, with the reason stated.
  \item Volatility is annualised with the correct factor and compared against a fiat baseline.
  \item Maximum drawdown is reported alongside the gain required to recover from it.
  \item Each of the three functions receives a separate verdict supported by a specific measurement.
  \item Synthetic data is identified as synthetic wherever a number is quoted.
\end{itemize}
\end{verification}

\begin{labdeliverable}
Submit the volatility and drawdown output, a completed three-function table with a verdict and supporting measurement for each row, and a two-page assessment: which function the asset satisfies best, which it fails, whether the failures are inherent to the design or fixable, and what would have to change for your verdict to reverse.
\end{labdeliverable}
